\section{Finite spatial radii under coexistence}\label{sec:radius}

The regeneration construction controls the advance of a slab in its chosen
direction.  We now control its full spatial extent, assuming that escape in
both directions has positive probability.  The idea is to turn a very large
slab into a new record for a slightly tilted linear functional.  From that
record a fresh path escaping in the opposite direction can preserve the
current height maximum forever.  If mean spatial radii were infinite, this
would give a fixed positive probability that the final maximum is finite
but arbitrarily large, which is impossible.

\begin{proposition}[Spatial radius under coexistence]\label{prop:radius}
Let the environment consist of iid strictly positive transition rows,
and let
\(v\in\RR^d\setminus\{0\}\).  Suppose
\(P_0(A_v)>0\) and \(P_0(A_{-v})>0\).  For each sign
\(\sigma\in\{+,-\}\), let \(\nu_{\sigma v}\) be the slab law from
Lemma~\ref{lem:slabs}.  If \(R\) is the maximum Euclidean distance from
a slab's initial site, including its terminal site, then
\[
 \int R\,d\nu_v+\int R\,d\nu_{-v}<\infty.
\]
\end{proposition}

This is a spatial moment statement.  It asserts no moment bound on the
number of steps in a slab.

We first record a maximal inequality that will control all initial portions
of a sequence of slabs at once. It is a consequence of Hopf's maximal
ergodic theorem; see Garsia's proof~\cite{Garsia1965}. The interval argument
below gives the precise stationary-sequence form that we use.

\begin{lemma}[Initial averages]\label{lem:initial-averages}
Let \((Z_j)_{j\ge1}\) be a stationary sequence of nonnegative integrable
random variables.  For every \(a>0\) and positive integer \(n\),
\[
 P\left(\max_{1\le k\le n}\frac1k\sum_{j=1}^kZ_j>a\right)
 \le \frac{EZ_1}{a}.
\]
\end{lemma}

\begin{proof}
Call an integer \(s\) a bad start if an interval beginning at \(s\), of
length at most \(n\), has sum greater than \(a\) times its length.
Among starts \(1,\ldots,M\), take the leftmost bad start, choose a
witness interval, skip all its indices, and repeat.  The chosen intervals
are disjoint, cover every bad start, and lie in \(\{1,\ldots,M+n\}\).
Their total length is at least the number of bad starts.  Therefore
\[
 a\,\#\{s\le M:s\text{ is bad}\}
 \le \sum_{j=1}^{M+n}Z_j.
\]
Taking expectations, using stationarity, and letting \(M\to\infty\)
proves the assertion.
\end{proof}

\begin{proof}[Proof of Proposition~\ref{prop:radius}]
The construction has two parts.  A displacement allowance \(a\) and a
number \(N\) of slabs will be chosen so that most initial portions have
total radius at most \(a\) times their length, but some slab has a much
larger excursion with probability bounded below.  That excursion supplies
a record for a small tilt of the opposite height direction.  The
continuation then avoids the old prefix in two successive regions: it
stays beyond the tilted record for a fixed multiple of \(N\) slabs, and
after their final true cut it stays at negative height.  The estimates
below arrange these two separations with a continuation probability
independent of \(a\) and \(N\).

Rescale \(v\) so that \(\|v\|_\infty=1\).  For either sign, under
\(\widehat P_{\sigma v}\) let \((L_i,R_i)_{i\ge1}\) denote the widths
and radii of its iid slabs, and put \(H_k=\sum_{i=1}^kL_i\), with
\(H_0=0\).  Write \(E_\sigma\) for expectation in this slab sequence.
Lemma~\ref{lem:mean-width} gives
\(0<E_\sigma L<\infty\).  A point in slab \(k+1\) has signed height
at least \(H_k\), and the displacement of any point through slab
\(k\) has norm at most \(\sum_{j=1}^kR_j\).

\paragraph{Scales with controlled initial portions.}
Suppose, to obtain a contradiction, that
\[
 I(t):=\sum_{\sigma\in\{+,-\}}E_\sigma\min(R,t)
 \longrightarrow\infty.
\]
The function \(I\) is increasing and concave, \(I(0)=0\), and
\(I(t)=o(t)\).  The last assertion follows by bounded convergence from
\(R<\infty\) almost surely.

By the two width strong laws we may fix \(c>0\) and finite
\(C,C_0\ge0\) such that, in each sign,
\begin{equation}\label{eq:radius-height-bounds}
 \widehat P_{\sigma v}
 \bigl(ck-C_0\le H_k\le Ck+C_0\text{ for every }k\ge0\bigr)>.95.
\end{equation}
Indeed choose \(c\) below both means and \(C\) above both means, then
choose \(C_0\) to bound the almost surely finite maximal deviations
with the displayed probability.  Fix, in this order, an integer \(C_1\),
numbers \(b,B\), and \(\eta>0\), so that
\begin{equation}\label{eq:radius-constants}
 C_1c>C+4,\qquad b>1/c,\qquad
 \frac{B-1}{\sqrt d}-b(C+2)>2,\qquad 4C_1\eta<.15.
\end{equation}
The multiplier \(C_1\) gives the opposite continuation enough slabs to
descend below height zero.  The coefficient \(b\) makes its height
descent dominate its lateral displacement in the tilted functional;
\(B\) ensures that a large excursion creates a new tilted record.
Finally, \(\eta\) controls the probability that an initial portion exceeds
its displacement allowance.  These constants remain fixed throughout
the proof.

For large \(a\), let \(N=N(a)\) be the smallest positive integer
such that \(I(aN)\ge\eta a\).  Such an integer exists.  Since
\(I(am)/a\to0\) for every fixed \(m\), we have \(N(a)\to\infty\).
For sufficiently large \(a\), minimality and concavity give
\begin{equation}\label{eq:radius-scale}
 \eta a\le I(aN)
 \le \frac{N}{N-1}I(a(N-1))<2\eta a.
\end{equation}

In either sign and for \(m\ge0\), define \(\mathcal R_m\) to be the
event that, for every \(k\le m\),
\[
 ck-C_0\le H_k\le Ck+C_0,
 \qquad \sum_{j=1}^kR_j\le ak.
\]
We claim that
\begin{equation}\label{eq:radius-good-prefix}
 \widehat P_{\sigma v}(\mathcal R_{C_1N})>.8.
\end{equation}
To see this, put \(n=C_1N\), \(A=an\), and truncate each radius to
\(\widehat R_j=\min(R_j,A)\).  Concavity and
\eqref{eq:radius-scale} give
\[
 E_\sigma\widehat R\le I(A)\le C_1 I(aN)<2C_1\eta a.
\]
Lemma~\ref{lem:initial-averages} bounds by \(2C_1\eta\) the
probability that any of the first \(n\) averages of the truncated
radii exceeds \(a\).  Also
\[
 \widehat P_{\sigma v}(\exists j\le n:R_j>A)
 \le n\widehat P_{\sigma v}(R>A)
 \le \frac{nE_\sigma\min(R,A)}{A}
 \le2C_1\eta.
\]
Intersecting these estimates with \eqref{eq:radius-height-bounds}
proves \eqref{eq:radius-good-prefix}.  In particular,
\(\widehat P_{\sigma v}(\mathcal R_{i-1})>.8\) for \(i\le N\),
and \(\mathcal R_{i-1}\) depends only on the first \(i-1\) slabs.

\paragraph{A large slab after a controlled initial portion.}
Choose a lower index \(i_0\) so large that
\begin{equation}\label{eq:radius-width-tail}
 \sum_\sigma\sum_{i\ge i_0}\widehat P_{\sigma v}(L>i)
 <\frac{\eta}{4B}.
\end{equation}
This is possible by integrability of the widths.  Increase \(i_0\)
if necessary so that, for every \(i\ge i_0\),
\begin{equation}\label{eq:radius-index-margins}
 c(i-1)-C_0\ge ci/2,
 \qquad C(i-1)+C_0+i\le(C+2)i.
\end{equation}
After fixing \(i_0\), take \(a\) sufficiently large that \(N\ge i_0\)
and \(I(Bai_0)<\eta a/4\).

For \(i_0\le i\le N\), call
\[
 G_i=\{R_i>Bai,\ L_i\le i\}
\]
the large-slab event, and put
\(\lambda_\sigma=\sum_{i=i_0}^N\widehat P_{\sigma v}(G_i)\).
Integrating the decreasing radius tails gives
\begin{align*}
 \sum_\sigma\sum_{i=i_0}^N\widehat P_{\sigma v}(R>Bai)
 &\ge \frac{I(Ba(N+1))-I(Bai_0)}{Ba},\\
 \sum_\sigma\sum_{i=1}^N\widehat P_{\sigma v}(R>Bai)
 &\le\frac{I(BaN)}{Ba}\le2\eta.
\end{align*}
For the last bound, use \(I(BaN)\le B I(aN)\).  For the lower
bound, use \(I(Ba(N+1))\ge I(aN)\ge\eta a\), then subtract
\eqref{eq:radius-width-tail}.  The result is
\[
 \frac{\eta}{2B}\le\lambda_++\lambda_-\le2\eta.
\]
Thus one sign, which may depend on the scale, satisfies
\(\lambda_\sigma\ge\eta/(4B)\).

In this sign consider
\[
 Z=\sum_{i=i_0}^N\1_{\mathcal R_{i-1}}\1_{G_i}.
\]
Independence of the \(i\)-th slab from its preceding slabs gives
\(E_\sigma Z\ge.8\lambda_\sigma\).  The summands of \(Z\) need
not be independent, but \(Z\le\sum_i\1_{G_i}\), whose summands
are independent.  Hence
\[
 E_\sigma Z^2\le\lambda_\sigma+\lambda_\sigma^2,
 \qquad
 \widehat P_{\sigma v}(Z>0)
 \ge\frac{.64\lambda_\sigma}{1+\lambda_\sigma}
 \ge\frac{.64\eta}{4B(1+2\eta)}=:\delta>0.
\]
The constant \(\delta\) is independent of \(i_0\) and \(a\).

We have therefore found, with a probability independent of the scale,
a slab whose spatial excursion is much larger than the accumulated
displacement before it, while its height width is at most its index.
The next step converts that excursion into a stopping time from which
an opposite-going path can use fresh environment rows.

\paragraph{A tilted record.}
Write \(h(x)=\sigma v\cdot x\) for the sign just selected.  On
\(\mathcal R_{i-1}\cap G_i\), the initial site of slab \(i\) has
norm at most \(a(i-1)\).  Some point of the slab consequently has
norm greater than \((B-1)ai\).  At that point some signed coordinate
\(u\in\{\pm e_1,\ldots,\pm e_d\}\) satisfies
\[
 u\cdot X>\frac{B-1}{\sqrt d}ai.
\]
Throughout this slab its running \(h\)-maximum is at least
\(H_{i-1}\ge ci/2\), and at most
\(H_i\le(C+2)i\), by \eqref{eq:radius-index-margins}.

Define the linear functional
\[
 Y(x)=u\cdot x-ba h(x).
\]
Before slab \(i\), nonnegative height and \(\mathcal R_{i-1}\)
give \(Y(X)\le a(i-1)\).  At the selected point,
\[
 Y(X)>\left(\frac{B-1}{\sqrt d}-b(C+2)\right)ai>2ai.
\]
There is thus a strict \(Y\)-record while the running \(h\)-maximum
lies in the deterministic window
\begin{equation}\label{eq:radius-window}
 W=[ci_0/2,(C+2)N].
\end{equation}

For each fixed choice of \(\sigma,u,i_0,a\), let \(T\) be the first
positive integer time such that the following three conditions hold:
the entire prefix has nonnegative \(h\)-height; \(Y(X_T)\) is a
strict record; and the running \(h\)-maximum belongs to \(W\).
These are conditions on the observed prefix, so \(T\) is a stopping
time for \((\cF_n)\).  It need not be a true regeneration time.
The preceding construction and a union over the \(2d\) signed
coordinates show that some \(u\) satisfies
\[
 \widehat P_{\sigma v}(T<\infty)\ge\frac{\delta}{2d}.
\]
Set \(p_*:=\min(p_v,p_{-v})>0\).  Removing the conditioning yields
\begin{equation}\label{eq:radius-launch-probability}
 P_0(T<\infty)\ge\frac{p_*\delta}{2d}.
\end{equation}
The site \(X_T\) has not appeared in its earlier prefix, since it is
a strict \(Y\)-record.

\paragraph{Preserving the maximum with a fresh continuation.}
Since \(\|v\|_\infty=1\), there is a nearest-neighbor step \(e\)
with \(h(e)=-1\).  From \(X_T\) force \(M\) repetitions of this
step, where the fixed integer \(M\) will be chosen shortly.  Each
step increases \(Y\) by at least \(ba-1\).  For all sufficiently
large \(a\), its successive sites have strictly increasing
\(Y\)-coordinates, starting at the strict record \(X_T\).
Consequently every scripted departure site is distinct and unused by
the prefix.  Conditional on any fixed prefix realizing \(T\), the
script has the exact raw probability
\[
 \bigl(\EE[\omega(0,e)]\bigr)^M\ge\alpha^M,
 \qquad \alpha:=\min_{e'\in U}\EE[\omega(0,e')]>0.
\]
Here independence is across departure rows, not across the entries in
a single row.  The constant \(\alpha\) is positive because \(U\) is
finite and every transition entry is positive almost surely.  It is
not a lower bound on the transition entries in individual environments.

At the script endpoint consider a translated opposite-going template
with law \(\widehat P_{-\sigma v}\), and require
\(\mathcal R_{C_1N}\) for its first \(C_1N\) true-cut slabs.
This requirement includes existence of those cuts and concerns the
entire template; they are not declared to be stopping times.  It has
conditional-template probability greater than \(.8\).  At any point
of its slab \(k+1\), where \(0\le k<C_1N\), the coordinate
displacement in direction \(u\) is at least \(-a(k+1)\), while
the displacement in the original height \(h\) is at most
\(-ck+C_0\).  Its total \(Y\)-increment from the launching record
is consequently at least
\begin{equation}\label{eq:radius-template-separation}
 \begin{split}
 M(ba-1)-a(k+1)+ba(ck-C_0)
 &=a\bigl(Mb-1-bC_0+(bc-1)k\bigr)-M.
 \end{split}
\end{equation}
Choose \(M\) once and for all so that \(Mb>1+bC_0+1\).
Since \(bc>1\), the expression in
\eqref{eq:radius-template-separation} is strictly positive for all
\(k\) in question and all sufficiently large \(a\).  Thus these
template sites avoid the entire original prefix, whose \(Y\)-values
are at most the launching record.

There are two additional separations to check.  First, the whole
opposite template stays at or below its own initial \(h\)-height,
whereas every script departure lies strictly above that height.
It therefore avoids all script departures, including in its infinite
tail.  Second, after \(C_1N\) template slabs its absolute
\(h\)-height is at most
\[
 (C+2)N-M-cC_1N+C_0<0
\]
for large \(N\).  This endpoint is a true opposite cut, so its
remaining tail stays below zero.  That tail avoids the original
prefix, all of whose heights are nonnegative.  No assumption that
the script itself stays above zero is needed.

These inequalities make the use of fresh rows exact.  Condition on a
fixed finite prefix realizing \(T\), and then on the forced script.
The conditional environment differs from the original iid law only
at departure sites of these two finite words.  Under an independent
raw law, the probability of
\(D_{-\sigma v}\cap\mathcal R_{C_1N}\) is at least
\(.8p_{-\sigma v}\).  Every path in that event avoids all the old
departure sites after translation to the script endpoint, by the
preceding separations.  Lemma~\ref{lem:departure-transfer} therefore
identifies its probability with the corresponding fresh-row
probability after the prefix and script.  More explicitly, kill both
the continuation and an independent raw chain on arrival at the same
set of old departure sites.  The finite-word identity gives equality
of their cylinder probabilities, hence of their killed path laws.
Now restrict to \(D_{-\sigma v}\cap\mathcal R_{C_1N}\), whose paths
never hit that set.  Repeated sites within the template are retained
in its raw path weights.

It follows that, from every qualifying prefix, the script followed by
the required continuation has conditional probability at least
\(.8p_*\alpha^M\).  This argument uses finite-prefix conditioning
and departure-row factorization, not a Markov property at a
future-dependent true cut.  The script decreases \(h\), and the
template never exceeds its own starting height.  Hence the final
supremum of \(h\) on this continuation is its running supremum at
\(T\), which belongs to \(W\).

\paragraph{The finite final maximum.}
Together with \eqref{eq:radius-launch-probability}, this proves that
for every sufficiently large \(i_0\), after choosing \(a\) large
enough, some sign \(\sigma\) satisfies
\begin{equation}\label{eq:radius-final-maximum}
 P_0\left(\sup_{n\ge0}\sigma v\cdot X_n
       \in[ci_0/2,(C+2)N(a)]\right)
 \ge \rho,
 \qquad
 \rho:=\frac{.8p_*^2\alpha^M\delta}{2d}>0.
\end{equation}
The constant \(\rho\) does not depend on the scale.  For each sign let
\(M_\sigma=\sup_{n\ge0}\sigma v\cdot X_n\), allowing the value
\(+\infty\).  Equation~\eqref{eq:radius-final-maximum} implies
\[
 \sum_{\sigma\in\{+,-\}}
 P_0\bigl(ci_0/2\le M_\sigma<\infty\bigr)\ge\rho
\]
for every sufficiently large \(i_0\).  Both terms on the left tend to
zero as \(i_0\to\infty\), a contradiction.  Thus \(I(t)\) is bounded, and
monotone convergence proves Proposition~\ref{prop:radius}.
\end{proof}

Proposition~\ref{prop:radius} supplies integrable slab displacements as well as
integrable errors between slab endpoints.  Section~\ref{sec:arrangement} uses these
spatial facts to compare the two limiting directions and then to arrange
the paths along a common coordinate height.
